\documentclass{article}
\usepackage{amsmath}
\usepackage{amsthm}
\usepackage{amsfonts}
\usepackage{amssymb}
\usepackage{bbold}
\usepackage{dsfont}
\usepackage{graphicx}
\usepackage{cleveref}
\usepackage{bm}
\usepackage{booktabs}
\usepackage{siunitx}
\usepackage{xcolor}
\usepackage[utf8]{inputenc}
\usepackage[english]{babel}
\newtheorem{definition}{Definition}
\newtheorem{example}{Example}
\newcommand{\exP}[1]{\mathrm{e}^{#1}}
\newcommand{\dif}{\mathrm{d}}
\renewcommand{\vec}[1]{\mathbf{#1}}
\newcommand{\pdagger}{{\phantom{\dagger}}}
\newcommand{\state}[1]{\vert #1\rangle}
\newcommand{\cstate}[1]{\langle #1\vert}
\newcommand{\tr}[1]{\mathrm{tr}\left\{#1\right\}}
\newcommand{\trs}[2]{\, \mathrm{tr}_{#1}\{#2\}}
\newcommand{\expc}[3]{\langle #1\vert #2 \vert #3\rangle}
\usepackage[doi=false,url=false,isbn=false,backend=bibtex,style=nature,natbib=true,sorting=none]{biblatex}
\bibliography{bib.bib}
\usepackage{geometry}
\geometry{
	a4paper,
	total={170mm,257mm},
	left=20mm,
	top=20mm,
}
\begin{document}
\title{Overview of collected data and outlook}
\author{Rudolf Smorka}
\maketitle
\section{Comparison between TTHEOM and MFHEOM (report from 02.03.23)}
Done:
\begin{itemize}
	\item Two-site spinless Peierls model: 
	\begin{itemize}
		\item MFHEOM and TTHEOM currents agree well for $t_1=\SI{4}{\milli\electronvolt}$ and $\omega_0 = \SI{2}{\milli\electronvolt},\SI{50}{\milli\electronvolt}, \SI{100}{\milli\electronvolt}$
		\item Disagreement for larger $t_1$: MFHEOM does not capture reduction of current with increasing $t_1$
	\end{itemize}
	\item Two-site spinfull Peierls and three-site Fransson models:
	\begin{itemize}
		\item Not able to obtain numerically converged results for $\omega_0=\SI{50}{\milli\electronvolt}$ and $t_1=\SI{20}{\milli\electronvolt}$ even with $r=140$, deviation of norm from unity is $\sim 0.002$ after $\SI{250}{\femto\second}$
	\end{itemize}
\end{itemize}
To be done:
\begin{itemize}
	\item Agreed with Raffaele to use the tAMEn solver, implementation in other programming language (?). Could potentially take weeks/months to have a working code if implemented by myself. 
\end{itemize}
\section{Analysis of transport in spinless two-site Peierls and comparison between analytical results and MFHEOM data (report from 10.03.23)}
Done:
\begin{itemize}
	\item Single-particle eigenenergies and eigenstates of the uncoupled and coupled problem
	\item Potential of mean force, and quantitative comparison between its minimum and vibrational position
	\item Qualitative explanation of current-voltage characteristics based on vibrational dynamics and how it affects single-particle eigenstates
\end{itemize}
To be done:
\begin{itemize}
	\item Extend analysis to spinfull two-site model and three-site Fransson model
	\item Compute I-V for system sizes $2\times 4, 3\times 4,4\times 4,5\times 4$. Spin polarization is expected to be larger for bigger systems. Question is if there is any spin polarization within Ehrenfest approach.
\end{itemize}
\section{Analysis of uniqueness of steady-state MFHEOM (report from 25.03.23)}
Done:
\begin{itemize}
	\item For parameters used by Fransson, the steady-state in Fransson and Peierls models do not depend on initial condition
	\item For $t_1=t_0=\SI{40}{\milli\electronvolt}$ even for two-site spinless molecule, steady-state depends on initial condition, model displays bistability. 
\end{itemize}
\end{document}




